Un diapason de \qty{\v{f}}{Hz} vibre au-dessus d'un long tube vertical partiellement rempli d'eau. On abaisse lentement le niveau de l'eau : on entend une résonance forte pour une certaine hauteur de la colonne d'air, puis la résonance suivante lorsque le niveau de l'eau est descendu de \qty{\v{dL}}{cm} de plus.

\begin{enumerate}[label=\alph*)]
    \item Expliquez pourquoi la distance entre deux résonances successives vaut une demi-longueur d'onde, puis déterminez la longueur d'onde du son.
    \item Déterminez la vitesse du son dans le tube.
    \item À l'aide de $v = \qty{331}{m/s}\sqrt{T/\qty{273}{K}}$, déterminez la température de l'air dans le tube, en degrés Celsius.
    \item En négligeant la correction d'extrémité, quelle est la plus petite longueur de colonne d'air qui donne une résonance ?
\end{enumerate}
